\section{Reweighting, WHAM}
I am doing enhanced sampling simulations of molecular systems.  There are a number of scenarios:
\begin{enumerate}
    \item Simgle unbiased ensemble
    \[
        \langle O\rangle = \int\myd{s}p(s)\,O(s)
        =\lim_{N\rightarrow\infty} \frac{1}{N}\sum_{k=1}^N O(s_k), ~s_k\sim p
    \]
    Unbiased distribution
    \[
    p(s) = \frac{\mye^{-\beta F(s)}}{Z}
    \]
    Biased distribution with time-dependent bias potential:
    \[
        p_t(s) = \frac{\mye^{-\beta [F(s)+ w_t(s)}}{Z_t}
    \]
    Reweighting factor for Monte Carlo sampling using biased simulation data:
    \[
        W_t(s_t) = \mye^{\beta[w_t(s_t)-c_t]} ,~~~\mye^{-\beta c_t} = \frac{Z_t}{Z}
    \]
    Estimating factors involving $c_t$ using biased simulation data
    \[
        \mye^{-\beta c_t} \approx \frac{\sum_{t'=1}^T W_{t'}(s_{t'})\mye^{-\beta w_t(s_{t'})}}{\sum_{t'=1}^T W_{t'}(s_{t'})}
    \]
    where the weights themselves depend on $c_t$.
\end{enumerate}

\subsection{Connection Between Reweighting and WHAM}

The key point is that there are two distinct kinds of weights:

\begin{enumerate}
    \item the \emph{single-ensemble reweighting factor}, and
    \item the \emph{multi-ensemble WHAM combination weights}.
\end{enumerate}

These serve different purposes.

\subsection{Single-Ensemble Reweighting}

Suppose the unbiased distribution is

\[
p(s)=\frac{e^{-\beta F(s)}}{Z}.
\]

A biased ensemble with time-dependent bias potential \(w_t(s)\) samples

\[
p_t(s)=\frac{e^{-\beta[F(s)+w_t(s)]}}{Z_t}.
\]

Define the reweighting factor

\[
W_t(s)=e^{\beta[w_t(s)-c_t]},
\]

with

\[
e^{-\beta c_t}=\frac{Z_t}{Z}.
\]

Then

\[
p(s)=p_t(s)e^{\beta[w_t(s)-c_t]}.
\]

Thus, the factor

\[
e^{\beta[w_t(s)-c_t]}
\]

removes the umbrella bias and converts samples from the biased ensemble into unbiased samples.

This is exactly the same structure as Roux Eq.~(7),

\[
[(p(\xi))]^{\mathrm{unbiased}}_{(i)}
=
e^{+\beta w_i(\xi)}
(p(\xi))_{(i)}
e^{-\beta F_i}.
\]

The correspondence is

\[
F_i \leftrightarrow c_i.
\]

Indeed, Roux defines

\[
e^{-\beta F_i}
=
\left\langle e^{-\beta w_i}\right\rangle
=
\frac{Z_i}{Z},
\]

which is identical to the normalization ratio above.

At this stage, we are still dealing with a \emph{single} ensemble.

\subsection{WHAM Introduces a Second Weighting Layer}

Suppose we now have many umbrella windows \(i=1,\dots,K\).

Each window produces an unbiased estimator

\[
\hat p_i(s)
=
p_i(s)e^{\beta[w_i(s)-F_i]}.
\]

If sampling were infinite, all \(\hat p_i(s)\) would agree exactly.

However, with finite sampling:

\begin{itemize}
    \item some windows are noisy,
    \item some regions are poorly sampled,
    \item overlap between windows may be uneven.
\end{itemize}

WHAM asks:

\[
\textit{How should all these estimators be combined optimally?}
\]

\subsection{The WHAM Equation}

Roux Eq.~(6) gives

\[
p(s)
=
\frac{
\displaystyle
\sum_i n_i p_i(s)
}{
\displaystyle
\sum_j n_j e^{-\beta[w_j(s)-F_j]}
}.
\]

This can be rewritten as

\[
p(s)
=
\sum_i a_i(s)\,\hat p_i(s),
\]

where

\[
a_i(s)
=
\frac{
n_i e^{-\beta[w_i(s)-F_i]}
}{
\displaystyle
\sum_j n_j e^{-\beta[w_j(s)-F_j]}
}.
\]

The factors \(a_i(s)\) are the WHAM combination weights.

Thus there are \emph{two} layers:

\begin{enumerate}
    \item Bias removal:
    \[
    e^{\beta[w_i(s)-F_i]},
    \]
    which converts a biased distribution into an unbiased one.

    \item Statistical combination:
    \[
    a_i(s),
    \]
    which combines all windows into a minimum-variance estimate.
\end{enumerate}

\subsection{Interpretation of the WHAM Weights}

Suppose \(s\) lies near the center of umbrella \(i\). Then

\[
w_i(s)\approx 0,
\]

so umbrella \(i\) samples that region efficiently.

For another umbrella \(j\) far away,

\[
w_j(s)\gg 1,
\]

and that window contributes very little statistical information there.

Consequently,

\[
a_i(s)
\propto
n_i e^{-\beta[w_i(s)-F_i]}
\]

gives large weight to windows that sample \(s\) well.

Thus:

\begin{itemize}
    \item the reweighting factor removes the bias,
    \item the WHAM factor determines statistical confidence.
\end{itemize}

These are conceptually distinct.

\subsection{Relation to Multi-Ensemble Monte Carlo}

WHAM is fundamentally a multi-ensemble importance sampling method.

Each umbrella window samples

\[
\pi_i(x)
\propto
e^{-\beta[U(x)+w_i(x)]}.
\]

The unknown constants \(F_i\) are relative free energies,

\[
F_i
=
-k_B T \ln \frac{Z_i}{Z}.
\]

WHAM determines the \(F_i\) self-consistently so that all windows agree on the same underlying unbiased distribution.

This is closely related to:

\begin{itemize}
    \item Bennett acceptance ratio (BAR),
    \item Ferrenberg--Swendsen reweighting,
    \item MBAR (multistate Bennett acceptance ratio).
\end{itemize}

In fact:

\[
\begin{aligned}
\text{BAR} &\rightarrow \text{two ensembles}, \\
\text{WHAM} &\rightarrow \text{histogram formulation}, \\
\text{MBAR} &\rightarrow \text{continuous/unbinned generalization}.
\end{aligned}
\]

\subsection{Answer to the Original Question}

Yes.

Roux Eq.~(7) gives the unbiased estimator from a single window.

WHAM then introduces an additional weighting factor to combine all windows optimally.

Schematically:

\[
\text{biased histogram}
\quad
\overset{\text{Eq.~7}}{\longrightarrow}
\quad
\text{unbiased local estimate}
\quad
\overset{\text{WHAM}}{\longrightarrow}
\quad
\text{global optimal estimate}.
\]

Thus the WHAM weights are an additional layer on top of the single-window reweighting factors.
